% Generated draft: requires mathematical and source review.
\documentclass[11pt]{article}
\usepackage[margin=25mm]{geometry}
\usepackage{fontspec}
\setmainfont{Latin Modern Roman}
\usepackage{amsmath,amssymb,mathtools,unicode-math}
\setmathfont{Latin Modern Math}
\usepackage{enumitem,needspace,xurl,hyperref}
\setlength{\parindent}{0pt}
\setlength{\parskip}{5pt}
\setlength{\emergencystretch}{3em}
\newcommand{\N}{\mathbb{N}}
\newcommand{\Z}{\mathbb{Z}}
\newcommand{\Q}{\mathbb{Q}}
\newcommand{\R}{\mathbb{R}}
\newcommand{\C}{\mathbb{C}}
\newcommand{\F}{\mathbb{F}}
\newcommand{\A}{\mathbb{A}}
\newcommand{\PP}{\mathbb{P}}
\newcommand{\E}{\mathbb{E}}
\newcommand{\eps}{\varepsilon}
\newcommand{\dd}{\,\mathrm{d}}
\begin{document}
\section*{UnsolvedMath ID: 20000818}
\textbf{Source problem code:} catalogue entry
\subsection{Definitions and mathematical statement}
The source fixes an integral surface \(S\subset\PP^3\) of degree \(d\), with a double curve \(D\) of degree \(e\) and singularities described only as ``triple points, pinch point, \(\ldots\).'' Its conjecture asks for an integer \(n_0\) such that every smooth curve \(C\subset S\) of degree \(n\ge n_0\) has the following property: a general curve \(C'\) in its Hilbert-scheme component \(H_C\) is contained in a degree-\(d\) surface \(S'\) having a double curve \(D'\) of degree \(e\). Here general means on a dense open subset of the component, and \(H_C\) is a component of the Hilbert scheme with the Hilbert polynomial of \(C\). The source does not specify a ground field, the full allowed singularity class, the dependence of \(n_0\), or which component to choose if \([C]\) lies on several. The double curve is interpreted as belonging to \(S'\), rather than requiring \(C'\subset D'\).

The supplied partial result reports that \(n>d^2\) makes \(S\) the unique degree-\(d\) surface containing \(C\), and excludes containing surfaces of smaller degree. It calls \(d^2\) the optimal uniform threshold for this uniqueness assertion; this does not supply the conjectural threshold \(n_0\) for preservation of the double curve.

For the aid's local deformation criterion, work over an algebraically closed field of characteristic zero, as a setting in which dominance and tangent criteria can be interpreted without inseparability issues; this field choice is not explicit in the original source. Specify a smooth locally closed stratum \(B\) of degree-\(d\) surfaces with simultaneous normalization and flat conductor data, preserving the prescribed ordinary singularity type and double-curve degree \(e\). The existence and precise construction of such a stratum are additional data, not supplied by the original ellipsis. Let \(F\) define \(S\), and define
\[
 \rho:H^0(C,N_{C/\PP^3})\longrightarrow H^0(C,\mathcal O_C(d)),
 \qquad \xi\longmapsto\xi(F),
\]
\[
 r:T_{[S]}B\longrightarrow H^0(C,\mathcal O_C(d)),
 \qquad \eta\longmapsto\text{restriction to }C\text{ of the equation deformation }\eta.
\]
For the restricted flag space of pairs \((C',S')\) with \(S'\in B\), the aid reports that the image of its projection differential to the curve Hilbert scheme is exactly
\[
 \{\xi\in H^0(C,N_{C/\PP^3}):\rho(\xi)\in\operatorname{im}r\}.
\]
Assuming \(n>d^2\), smoothness of the curve Hilbert point and flag point, and the containment
\[
 \rho\bigl(H^0(C,N_{C/\PP^3})\bigr)\subseteq\operatorname{im}r,
\]
the aid reports that this flag projection is locally an open immersion. Thus the general curve on the corresponding Hilbert component lies on a unique surface in \(B\), giving the desired double-curve degree for this specified local setting. Its tangent obstruction is
\[
 Q=\frac{\rho(H^0(C,N_{C/\PP^3}))}
 {\rho(H^0(C,N_{C/\PP^3}))\cap\operatorname{im}r}.
\]
It reports that nonvanishing of \(Q\) at a general flag obstructs dominance, subject to the same smoothness and characteristic qualifications. The original conjecture remains broader: the aid does not prove the containment uniformly for all sufficiently large-degree curves or construct an appropriate stratum for every source surface.
\subsection{Short English statement}
Must sufficiently high-degree smooth curves on a surface with a double curve deform only to curves on surfaces with the same degree and double-curve degree? The supplied aid gives uniqueness of the containing surface above a degree threshold and a conditional tangent-space test for preserving its conductor data.
\subsection{Sources}
\begin{itemize}
\item Local UnsolvedMath registry, permanent ID \(20000818\), source code \texttt{catalogue entry}; \texttt{aim-arithmetic-geometry-notes.json}, \texttt{notes[63]}, tag \texttt{problem}.
\item AIM workshop ``Components of Hilbert schemes,'' Problem 17: \url{https://aimath.org/WWN/hilbertschemes/hilbertschemes.pdf}. Statement status \texttt{exact}; no PDF was accessed for this draft.
\item P. Ellia and R. Hartshorne, \emph{Smooth specializations of space curves: questions and examples}, in \emph{Commutative algebra and algebraic geometry (Ferrara)}, pp. 53--79, Lecture Notes in Pure and Applied Mathematics 206, Dekker, New York, 1999. The original remark attributes the conjecture to this work.
\item J. Brevik and F. Mordasini, \emph{Curves on a ruled cubic surface}, Collect. Math. 54 (2003), no. 3, pp. 269--281. The original remark reports the conjecture proved for \(d=3\).
\item Research attempt 2: \texttt{partial\_result}, status \texttt{partially\_solved}, classification \texttt{PARTIAL-PROGRESS}, rigor label \texttt{proved}, novelty confidence low. Created and updated 14 August 2026. The machine-generated aid is not peer reviewed; independently verify before citation.
\item Ulam AI / UnsolvedMath Contributors; CC BY 4.0.
\end{itemize}
\end{document}
